\documentclass[11pt]{article}


\usepackage[margin=1in]{geometry}
\usepackage[T1]{fontenc}
\usepackage[utf8]{inputenc}
\usepackage[expansion=false]{microtype}
\usepackage{booktabs}
\usepackage{array}
\usepackage{graphicx}
\usepackage{xcolor}
\usepackage{amsmath}
\usepackage{enumitem}
\usepackage[hidelinks]{hyperref}
\usepackage{xurl}
\usepackage{listings}

\lstset{basicstyle=\ttfamily\small, breaklines=true, frame=single, framesep=4pt}


\title{Benign Memory Contamination in LLM Agents:\\
A Taxonomy, Benchmark, and Contract-Level Mitigations}

\author{Eduardo Arana}

\date{July 2026}

\begin{document}

\maketitle

\begin{abstract}
Persistent memory makes LLM assistants persistent collaborators, but it also
creates an unmanaged input channel into future responses. We study \emph{benign
memory contamination}: degradation caused by ordinary memory-system behavior,
without an attacker, jailbreak, or poisoned record. We contribute a taxonomy
of six mechanisms: semantic drift, provenance collapse, scope bleed, temporal
staleness, recursive compounding, and summarization loss. We also introduce
CONTAM-Bench, an open, spec-driven benchmark evaluated here as a nine-scenario
pilot with a personalization-retention control, and a proposed mitigation architecture
based on scoped namespaces, provenance tags, raw fidelity, and decay by fact
class.

We report an ablation over seven configurations using one subject model and
TF-IDF retrieval, followed by five repetitions of each scenario--configuration
cell. On the targeted probes, the namespacing arm excluded cross-domain
distractors and the TTL arm excluded the stale fact in all five repetitions.
On the provenance probe, tags and raw fidelity each yielded five
machine-resolved clean outcomes, consistent with preserving source information
that was lost during summarization in this case. In a seeded-recursion case,
logged relevance-gate decisions discarded premise-denying context while
retaining a contaminated write-back. We call this observed failure
\emph{presupposition capture} and propose contradiction preservation as a
safeguard whose downstream efficacy remains untested. Of 350 scored rounds,
52 remain unresolved pending human review. These are within-scenario
observations, not estimates of real-world prevalence or model-independent
mitigation efficacy. The benchmark and
code are available at \url{https://github.com/arananet/contam-bench}.
\end{abstract}

\section{Introduction}
\label{sec:intro}

Memory transforms an LLM assistant from a stateless tool into a persistent
collaborator. It is also an unmanaged input channel into every future
response. The failure mode studied here requires no adversary: an
assistant's own retrieval pipeline injects semantically similar but
task-irrelevant context from past sessions, producing wrong assumptions,
provenance errors, and, when contaminated outputs are written back to the
store, compounding degradation over time. We call this \emph{benign
memory contamination} to distinguish it from adversarial memory poisoning,
where malicious records are deliberately injected~\cite{chen2024agentpoison,
dong2025minja}.

The problem is not hypothetical, and it is not confined to consumer
personalization. Recent survey work on harness engineering argues that the
realistic near-term path to recursive self-improvement runs through the
harness layer, including workflow design, tool orchestration, evaluation, and
context and memory lifecycle management, rather than direct
weight modification~\cite{weng2026harness}. That survey explicitly lists
memory and context degradation in long-horizon agent loops as an unresolved
structural challenge. This paper studies a bounded part of that problem:
it operationalizes six candidate mechanisms in hand-authored scenarios and
examines how selected memory controls affect retrieval and response outcomes.

A live illustration motivated this work. During the drafting of this very
paper, an LLM assistant with persistent memory retrieved the author's own
open-source specification proposal from a prior session and, unprompted,
elevated it to ``the agentic memory standard,'' a status-inflation and
provenance error produced by similarity-based cross-session retrieval,
exhibiting two of the taxonomy classes defined in
Section~\ref{sec:taxonomy} before the benchmark existed to name them.

\paragraph{Contributions.}
\begin{enumerate}[leftmargin=*, itemsep=2pt]
  \item A taxonomy of six benign contamination mechanisms, each defined by
        its causal mechanism and a falsifiable contamination criterion,
        plus one documented experimental class
        (Section~\ref{sec:taxonomy}).
  \item CONTAM-Bench, an open, spec-driven benchmark: hand-authored YAML
        scenario manifests, a pluggable memory store spanning five design
        axes, deterministic-first scoring with subject/judge model
        separation, and six metrics including a personalization-retention
        control (Section~\ref{sec:benchmark}).
      \item A proposed mitigation architecture expressed as declarative
      memory contracts, including scoping, provenance, gating, and decay,
      placed
        deliberately outside the mutable memory loop
        (Section~\ref{sec:architecture}).
      \item A seven-configuration pilot ablation with bundled-metadata limitations
        and a separately reported five-repetition audit on a Sonnet-class
        subject model
        (Section~\ref{sec:results}).
\end{enumerate}

\section{Related Work}
\label{sec:related}

The study draws on work on memory poisoning, non-adversarial memory risk,
retrieval noise, and memory-system evaluation. Its contribution is a small,
auditable mechanism-oriented test suite, not the discovery of memory risk
or a comprehensive comparison of existing memory systems.

\paragraph{Adversarial memory poisoning.}
AgentPoison demonstrated backdoor attacks that poison agent long-term memory
and RAG knowledge bases~\cite{chen2024agentpoison}; MINJA showed malicious
records can be injected purely through query-level
interaction~\cite{dong2025minja}; subsequent work has extended the attack
surface. This literature assumes an attacker and optimizes defenses
accordingly. Benign contamination has no attacker to defend against: the
threat model is the design of the memory system itself.

\paragraph{Measurement of non-adversarial memory risk.}
The closest neighbor is \emph{Remembering More, Risking
More}~\cite{altawaha2026remembering}, which measures memory-induced
violation rates exceeding no-memory baselines, increasing with exposure
length, and driven by cross-context contamination, identifying retrieval
scope and summarization level as amplifying design choices. Complementary
diagnosis work finds most memory failures stem from irrelevant retrieval
rather than memory construction. In the multi-agent setting, Wang et
al.~\cite{wang2026state} demonstrate \emph{memory laundering}: toxic
context compressed into memory summaries evades standard detectors while
still influencing downstream generations, and mitigation depends on
intervening before compression into persistent state. Their focus is
toxicity propagation; this paper's is task-correctness and attribution
degradation from memory operating as designed, but both establish that
contamination survives below detection thresholds. These papers establish
that the benign problem is real and measurable. Here we focus on connecting
an explicit scenario rubric to a selected memory control and a separately
inspectable retrieval trace; this is complementary to those measurements,
not evidence that the taxonomy exhausts the mechanisms they study.

\paragraph{Retrieval noise and context failure.}
The retrieval-layer half of this paper's phenomenon has a single-shot
ancestor: Cuconasu et al.\ show that irrelevant passages in RAG divide
into random ones, which do not degrade answer quality, and
\emph{distracting} ones, which do~\cite{cuconasu2024power}, and
follow-on work quantifies a passage's distracting effect as a graded
severity measure. This paper's contribution relative to that line is the
persistence layer and the verdict split: memory adds write-back, decay,
and provenance dynamics that single-shot retrieval does not have, and
the retrieval-layer/response-layer dissociation of
Section~\ref{sec:twolayer} is measured across that loop. The distracting
effect is a candidate severity proxy for the full benchmark, superior to
forbidden-pattern hit counts. Separately, practitioner taxonomies of
context failure, notably Breunig's context rot classes of poisoning,
distraction, confusion, and clash~\cite{breunig2025context}, overlap
with several classes here at the symptom level; the differentiator of
this taxonomy is that it is memory-scoped, mechanism-causal (each class
is defined by the design property that produces it), and bound to a
prescriptive control.

\paragraph{Memory systems and their evaluation.}
A rich systems literature builds agent memory capability: hierarchical
context management~\cite{packer2023memgpt}, forgetting-curve
decay~\cite{zhong2024memorybank}, temporal knowledge
graphs~\cite{rasmussen2025zep}, and production memory
stacks~\cite{chhikara2025mem0, xu2025amem}. Evaluation has followed:
MemBench~\cite{tan2025membench} measures memory effectiveness, and Zhou et
al.~\cite{zhou2026agentnative} systematically benchmark twelve memory
systems from a data-management perspective, decomposing memory into
storage, extraction, retrieval, and maintenance modules and finding that no
single architecture dominates. Closest in spirit to this paper's concern,
Xiong et al.~\cite{xiong2025experiencefollowing} empirically characterize
experience-following behavior, showing how past records steer agent
decisions without any adversary. These works decompose the \emph{system}
and measure capability, robustness, and cost; this paper decomposes the
\emph{failure mode} and measures contamination per class against
prescriptive design contracts.

\paragraph{Harness engineering.}
A recent survey organizes the design space of agent harnesses, including workflow
automation, filesystem-backed memory, sub-agent orchestration, and
self-improving harness loops such as ACE, Meta-Harness, Self-Harness, and
the Darwin G\"odel Machine~\cite{weng2026harness, zhang2025dgm}. The survey
treats memory hygiene as an open problem and predicts that harness
configuration and interface protocols will standardize industry-wide.

The positioning is therefore: surveys map the territory; poisoning papers
defend the border; measurement papers quantify the internal decay; this
paper names the decay's mechanisms and prescribes the architecture.

\section{A Taxonomy of Benign Contamination}
\label{sec:taxonomy}

Benign memory contamination is the degradation of assistant responses caused
by persistent memory operating as designed. Every class below arises from
ordinary usage patterns. The classes index onto the four stages of the
memory loop, which is also where their contracts bind: summarization loss
is a \emph{write-time} failure; semantic drift, scope bleed, and temporal
staleness are \emph{retrieve-time} failures; authority inflation is an
\emph{inject-time} dynamic; and recursive contamination is a
\emph{write-back} dynamic that compounds whichever of the others is
present. Recursion is thus not a peer of the static classes but a loop
over them, which is why it is tested both naturally (does the subject
generate the error?) and seeded (what happens when one exists?). Six classes are tested in the validation run; a
seventh is documented as experimental. Each class is defined by
(a)~its causal mechanism and (b)~a falsifiable contamination criterion used
verbatim by the benchmark's judge.

\subsection{Semantic drift (\texttt{semantic\_drift})}
Retrieval surfaces a memory because of surface-token overlap, not semantic
relevance, and the response is steered by the irrelevant memory.
\emph{Mechanism:} lexical or embedding similarity between unrelated topics
(``containers'' in Kubernetes vs.\ cooking; ``hosting'' a service vs.\ a
dinner). \emph{Contaminated when:} the response injects the unrelated
technical context into the answer.

\subsection{Provenance collapse (\texttt{provenance\_collapse})}
An assistant-sourced statement is later replayed as if it were a user-stated
fact or preference; the memory system dropped or never recorded who said it.
\emph{Mechanism:} untagged storage erases the user/assistant/tool
distinction; a suggestion becomes a preference on replay.
\emph{Contaminated when:} the response asserts the assistant's suggestion as
the user's own stated preference.

\subsection{Scope bleed (\texttt{scope\_bleed})}
Memory from one life domain (e.g.\ work) leaks into responses in another
domain (e.g.\ personal) on the strength of topical overlap.
\emph{Mechanism:} a global, non-namespaced store retrieves across domain
boundaries. \emph{Contaminated when:} work-specific details surface in the
personal answer (or the inverse).

\subsection{Temporal staleness (\texttt{temporal\_staleness})}
An outdated fact wins over, or blends with, its newer replacement because
the store has no recency or expiry mechanism. \emph{Mechanism:} no decay;
similarity ranking is blind to age; older, richer memories can outrank newer
corrections. \emph{Contaminated when:} the response acts on the stale fact.

\subsection{Recursive contamination (\texttt{recursive})}
A mildly contaminated response is written back to the store and compounds:
the system consumes its own error as ground truth in later rounds.
\emph{Mechanism:} write-back of assistant output; second-round retrieval
surfaces the first-round contamination alongside the original seeds.
\emph{Contaminated when:} the second-round output repeats or amplifies the
first-round error, measured by the compounding factor
(Section~\ref{sec:metrics}).

This class gains external support from two independent results. First, the
STOP experiments showed that recursive self-improvement scaffolds improved
outcomes with strong models but \emph{degraded} them with weak
ones~\cite{zelikman2023stop}: recursion amplifies whatever is underneath.
CONTAM-Bench accordingly measures compounding per model rather than
assuming a constant. Second, the harness literature warns that when a
program can edit its own operating layer, abstraction boundaries break and
safety controls must live outside the loop~\cite{weng2026harness}, an
argument this paper extends from code to memory in
Section~\ref{sec:architecture}.

\subsection{Summarization loss (\texttt{summarization\_loss})}
Storing a summary instead of the raw memory drops the disambiguating
detail; a later probe hinging on that detail gets a wrong or overconfident
answer. \emph{Mechanism:} lossy compression at write time; the detail needed
later was not predictable at summarization time. \emph{Contaminated when:}
the response asserts the wrong disambiguation or invents the missing detail.

\subsection{Authority inflation (experimental, untested)}
A tool- or assistant-sourced memory acquires unwarranted epistemic authority
through repetition: each retrieval-and-restatement cycle makes the fact look
more established (``as we've discussed several times\ldots''), even though
it traces to a single low-confidence origin. Distinguishing this cleanly
from provenance collapse (origin \emph{lost}) and recursive contamination
(error \emph{amplified}) requires multi-round scenarios with citation-count
probes that are out of scope for the validation test; the class is
documented in the benchmark taxonomy and excluded from all metrics.

\section{CONTAM-Bench}
\label{sec:benchmark}

CONTAM-Bench is a spec-driven benchmark: all scenarios are YAML manifests
validated against a published schema; the harness reads specs and never
hardcodes scenario logic; every run artifact, including raw prompts, injected
memory, responses, and judge verdicts, is persisted as JSON for audit. This
paper reports the scaled-down validation corpus. The base corpus is
\emph{eight} hand-authored scenarios: six contamination probes, one per
taxonomy class, plus two controls. The ablation of
Section~\ref{sec:results} adds a ninth, seeded-recursion scenario,
yielding nine scenarios against seven memory configurations. Hand-authoring
matters: a contamination pipeline cannot be validated with LLM-generated
scenarios without introducing contamination into its own ground truth.

\subsection{Scenarios}
Each contamination scenario seeds a memory store with hand-authored items
(each carrying content, source, age in days, domain, and fact class) and
then issues a probe query whose correct answer must not draw on the seeded
distractors. Expected outcomes are encoded as forbidden-content patterns.
Two controls complete the corpus: one where memory is genuinely relevant and
the correct behavior is to \emph{use} it, and one with an empty store, which
verifies the pipeline itself induces no artifacts.

\subsection{Memory design configurations}
The initial validation compared two bundled configurations across five design
axes (scope, storage fidelity, retrieval, provenance, decay). The v0.2
mechanism-isolation study adds five single-control arms to that comparison.

\begin{table}[h]
\centering
\small
\begin{tabular}{@{}lll@{}}
\toprule
Axis & \texttt{naive} & \texttt{governed} \\
\midrule
Scope & global store & namespaced by domain \\
Fidelity & summarized & raw \\
Retrieval & similarity top-$k$ ($k{=}4$) & top-$k$ + relevance gate \\
Provenance & untagged & tagged \texttt{[source: \ldots]} \\
Decay & none & TTL by fact class (365/90/14 d) \\
\bottomrule
\end{tabular}
\caption{Memory configurations under test. TTL classes: stable (365 days),
volatile (90), ephemeral (14); unclassified memories default to stable, the
recall-conservative and staleness-permissive choice. All thresholds are design
choices under test, not empirical constants; rationale is recorded in the
public spec.}
\label{tab:configs}
\end{table}

\subsection{Models and determinism}
The subject model under test is \texttt{claude-sonnet-4-6} at temperature~0;
the judge and relevance gate are \texttt{claude-haiku-4-5}. Subject and
judge are required to differ: a model never judges its own output. The
reported validation runs use TF-IDF cosine similarity, a documented
substitution because the public API exposes no embeddings endpoint. The
successor-study implementation additionally provides a local ONNX embedding
condition (\texttt{fastembed} with \texttt{BAAI/bge-small-en-v1.5}), reported
separately from TF-IDF; this paper reports no learned-embedding result.

\subsection{Verdict resolution and metrics}
\label{sec:metrics}

Scoring is deterministic-first. Each (scenario $\times$ configuration
$\times$ round) artifact receives exactly one verdict: a case-insensitive
regex pass over the scenario's forbidden-content patterns; a match yields
\texttt{contaminated}, otherwise a provisional \texttt{clean} that the LLM
judge may confirm. When deterministic and judge verdicts disagree, the
artifact is flagged \texttt{needs\_human\_review} and is never
auto-resolved; flagged artifacts are excluded from both numerator and
denominator of every rate, and every exclusion is listed in the report.

Six metrics are computed per configuration from persisted artifacts only;
no metric is ever invented, and an uncomputable metric is reported as null
with a machine-readable reason:

\begin{description}[leftmargin=*, itemsep=2pt]
  \item[Contamination rate.] Fraction of contamination probes judged
        contaminated. For the recursive scenario, the second-round verdict
        is used.
  \item[Provenance error rate.] The same rate restricted to provenance
        probes.
  \item[Staleness rate.] The same rate restricted to temporal-staleness
        probes. A retrieval assertion can establish that TTL removed a stale
        memory; a response-layer comparison remains null when its comparator
        is flagged for human review.
  \item[Compounding factor.] Ratio of second-round to first-round
        contamination severity for the recursive scenario, with severity
        measured as the count of distinct forbidden-content patterns
        matched. Values above~1 indicate amplification, 1 persistence,
        below~1 attenuation. A clean first round yields null: there is
        nothing to compound, which is the desired outcome.
  \item[Seeded-recursion rate.] Contamination rate for the separately
        reported seeded-recursion scenario. It is never pooled into the
        natural-probe contamination rate.
  \item[Personalization retention.] Fraction of control scenarios where the
      subject behaved correctly, using genuinely relevant memory and
        answering normally from an empty store without hallucinating
        remembered context. This control exists so that a design cannot
        ``solve'' contamination by never using memory.
\end{description}

\section{Ablation Study}
\label{sec:results}

Following external review of an earlier draft, the empirical section was
redesigned from a two-configuration validation into a single-control
ablation: the run reported here executes nine scenarios (the eight of
Section~\ref{sec:benchmark} plus a seeded-recursion scenario described
below) against seven configurations: \texttt{naive}, \texttt{governed},
and five arms that each add exactly one control to the naive
baseline.\footnote{Complete run artifacts are published at
\texttt{evidence/20260713T084130Z} in the benchmark repository, frozen at
tag \texttt{v0.2-ablation}. Every number in this section is recomputable
from those artifacts.} We characterize this study honestly: it is a
mechanism-isolation experiment on a small hand-authored corpus, designed
to validate the pipeline and attribute effects to individual controls, not as
a utility frontier. The latter requires the full corpus and is named as
successor work in Section~\ref{sec:discussion}. The run cost 132 API
calls (70 subject, 27 gate, 35 judge). The naive and governed
configuration hashes are unchanged from the prior validation run,
preserving comparability.

\begin{table}[h]
\centering
\small
\begin{tabular}{@{}lll@{}}
\toprule
Evidence release & Immutable tag & Backs in this paper \\
\midrule
v0.2 ablation & \texttt{v0.2-ablation} & Table~\ref{tab:matrix} single-run matrix \\
v0.3 repeated audit & \texttt{v0.3-repeated-ablation} & Tables~\ref{tab:repeated} and \ref{tab:retrieval} \\
v0.3.1 corrections & \texttt{v0.3.1-evidence-corrections} & call reconciliation and review queue \\
\bottomrule
\end{tabular}
\caption{Evidence lineage. Releases are not pooled: v0.2 is the frozen
single-run mechanism matrix; v0.3 is the repeated machine-resolved audit;
v0.3.1 is append-only accounting and adjudication metadata.}
\label{tab:lineage}
\end{table}

\paragraph{Seeded recursion (CB-VAL-009).} The natural recursion scenario
(CB-VAL-005) measures whether the subject \emph{generates} a round-1
error; with a strong subject it reliably does not, yielding a null
compounding factor. This is a model-strength result, not a mechanism test. The
seeded variant removes that dependence: the store is seeded directly with
an already-contaminated assistant-attributed write-back (``launch is
confirmed\ldots leadership has signed off,'' 2 days old) alongside the
user-sourced ground truth (``tentative Q4, nothing decided,'' 15 days
old), and the probe asks for the confirmed date. This measures write-back
propagation of the recursive mechanism, independent of subject-model
strength.

\subsection{The ablation matrix}

\begin{table}[h]
\centering
\small
\setlength{\tabcolsep}{4pt}
\begin{tabular}{@{}lccccccc@{}}
\toprule
Scenario (class) & \texttt{naive} & \texttt{ns} & \texttt{prov} & \texttt{ttl} & \texttt{gate} & \texttt{raw} & \texttt{gov} \\
\midrule
001 drift        & \(\circ\) & \(\circ\) & \(\circ\) & \(\circ\) & \(\circ\) & \(\circ\) & \(\circ\) \\
002 provenance   & \(\times\) & \(\times\) & \(\circ\) & \(\times\) & \(\times\) & \(\circ\) & \(\circ\) \\
003 scope bleed  & \(\times\) & \(\circ\) & \(\circ\) & \(\circ\) & \(\circ\) & \(\circ\) & \(\circ\) \\
004 staleness    & r & r & \(\circ\) & \(\circ\) & r & r & \(\circ\) \\
005 recursion (natural) & \(\circ\) & \(\circ\) & \(\circ\) & \(\circ\) & \(\circ\) & \(\circ\) & \(\circ\) \\
006 summarization & \(\circ\) & \(\circ\) & \(\circ\) & \(\circ\) & \(\circ\) & \(\circ\) & \(\circ\) \\
009 recursion (seeded) & r & r & \(\circ\) & r & \(\boldsymbol{\times}\) & r & r \\
\midrule
007 control (retention) & \(\circ\) & \(\circ\) & \(\circ\) & \(\circ\) & \(\circ\) & \(\circ\) & r \\
008 control (empty)     & \(\circ\) & \(\circ\) & \(\circ\) & \(\circ\) & \(\circ\) & \(\circ\) & \(\circ\) \\
\midrule
Contamination rate (natural probes) & 2/5 & 1/5 & \textbf{0/6} & 1/6 & 1/5 & 0/5 & 0/6 \\
Seeded recursion (009), separate & r & r & \textbf{0/1} & r & 1/1 & r & r \\
\bottomrule
\end{tabular}
\caption{Resolved verdicts per scenario and configuration
(\(\circ\)~clean, \(\times\)~contaminated, r~flagged
\texttt{needs\_human\_review} and excluded from rates per the metric
specification). Arms: \texttt{ns}~namespacing, \texttt{prov}~provenance
tags, \texttt{ttl}~decay, \texttt{gate}~relevance gate,
\texttt{raw}~raw fidelity; each is \texttt{naive} plus exactly that
control. Denominator policy (stated, not implied): the contamination rate
is computed over the six natural probes only; the seeded-recursion
scenario is reported as its own metric and never merged into the headline
rate. All rates show resolved counts; flagged artifacts are excluded from
numerator and denominator.}
\label{tab:matrix}
\end{table}

In the frozen v0.2 snapshot, ten artifacts were flagged for manual
inspection.
Most were quotation-or-negation cases: the subject \emph{quoted} a
forbidden claim while debunking it (``the memory says confirmed, but
nothing has actually been decided''), which fires the deterministic pattern
even when the response appears substantively clean. The accompanying v0.2
defect notes record an author's post-hoc assessment, not a blinded,
independent adjudication under a preregistered rubric. The v0.3 audit below
therefore replaces any inference from this small historical set with an
explicit repeated-run accounting of scorer disagreement.

\subsection{Repeated-evaluation audit and scoring defects}
\label{sec:repeatedaudit}

Table~\ref{tab:matrix} is the frozen v0.2 single-run ablation. We do not
pool it with the subsequent v0.3 repeated-evaluation audit: the latter ran
the same nine scenarios and seven configurations five times, producing 350
scored rounds, and is published separately at
\texttt{evidence/20260713T191740Z} under the tag in
Table~\ref{tab:lineage}. Its purpose is to quantify scorer ambiguity,
execution cost, and per-cell outcome distributions rather than replace the
matrix with an unqualified aggregate.

Fifty-two of the 350 rounds (14.9\%) required human review:
three provenance-attribution rounds had a
deterministic-clean/judge-contaminated disagreement; twenty staleness rounds
had a deterministic-contaminated/judge-clean disagreement; and twenty-nine
seeded-recursion rounds disagreed (twenty-eight in the latter direction, one
in the former). A candidate explanation for many disagreements is assertion-versus-mention:
the response quotes or rejects a stale or contaminated claim while the regex
detects its tokens. Conversely, a judge can label a provenance-sensitive
assertion contaminated when no forbidden string matches. Neither direction
establishes which scorer is correct without independent assessment. Per the
scoring contract, these rounds are excluded from machine-only rates;
agreement on the remaining rounds is not itself ground-truth validation.

Table~\ref{tab:repeated} reports repeated outcome distributions for the
load-bearing cells. Each entry is clean/contaminated/flagged across five
repetitions; non-flagged counts are machine-resolved only. The provenance
separation persists without flags: provenance tags, raw fidelity, and the
governed bundle are each 5/0/0, while naive is 0/4/1. The three provenance
disagreements occur only in \texttt{arm\_gate} (repetitions 1 and 5) and
\texttt{naive} (repetition 4), not in the tag, raw, or governed conditions.
Their artifact hashes are respectively \texttt{12e31cb370551aee},
\texttt{d4a2659218ff8d5d}, and \texttt{5e87b7133e0a98a5}.
This supports the provenance finding in the machine-resolved tier, but does
not substitute for the pending blinded human adjudications.

\begin{table}[p]
\centering
\scriptsize
\setlength{\tabcolsep}{3pt}
\begin{tabular}{@{}lcccc@{}}
\toprule
Configuration & 002 provenance & 003 scope & 004 staleness & 009 seeded recursion \\
\midrule
naive & 0/4/1 & 1/4/0 & 0/0/5 & 0/0/5 \\
namespace & 0/5/0 & 5/0/0 & 0/0/5 & 0/0/5 \\
provenance & 5/0/0 & 4/1/0 & 5/0/0 & 0/2/3 \\
ttl & 0/5/0 & 2/3/0 & 5/0/0 & 0/0/5 \\
gate & 0/3/2 & 5/0/0 & 0/0/5 & 0/4/1 \\
raw & 5/0/0 & 3/2/0 & 0/0/5 & 0/0/5 \\
governed & 5/0/0 & 5/0/0 & 5/0/0 & 0/0/5 \\
\bottomrule
\end{tabular}
\caption{Repeated v0.3 outcome distributions, reported as
clean/contaminated/flagged across five repetitions. Flagged means
\texttt{needs\_human\_review}; all other counts are machine-resolved.
Scenarios are not pooled into a population estimate.}
\label{tab:repeated}
\end{table}

The v0.3 evidence now includes a separately versioned queue for all 52
unresolved artifact rounds, but it contains no human verdicts. Therefore this
paper makes no human-consensus claim and does not silently adopt the judge
verdict. The rendered report now labels its machine-only and human-consensus
tables separately and identifies review rows by repetition and artifact hash.
The immutable original report still records the missing-adjudications state;
the append-only correction bundle and pending queue are identified in
Table~\ref{tab:lineage}.

The repeated audit used 660 API calls: 350 subject calls, 135 gate calls, and
175 judge calls. This exceeds the validation plan's stated ``few hundred
calls'' target. We report the overrun as an operational cost limitation; it
does not invalidate the persisted artifacts or their recorded results.

\subsection{The two-layer dissociation persists}
\label{sec:twolayer}

Retrieval-layer and response-layer outcomes differ on the semantic-drift
probe. Table~\ref{tab:retrieval} shows that namespacing, gating, and the
governed bundle excluded its distractors, whereas the other arms did not.
Response scoring nevertheless resolved clean across configurations. This
case illustrates why retrieval of an irrelevant memory and a contaminated
response must be measured separately. We therefore report three
non-interchangeable layers:
deterministic retrieval assertions for drift and scope; response verdicts
for attribution, staleness, and write-back propagation; and a utility layer
for summarization, which is \texttt{null} in this corpus because no utility
oracle has yet been declared. A clean response verdict is not a utility
measurement.

\begin{table}[h]
\centering
\small
\begin{tabular}{@{}lccccccc@{}}
\toprule
Retrieval assertion & naive & ns & prov & ttl & gate & raw & gov \\
\midrule
001 exclude drift distractor & 0/5 & 5/5 & 0/5 & 0/5 & 5/5 & 0/5 & 5/5 \\
003 exclude cross-domain distractor & 0/5 & 5/5 & 0/5 & 0/5 & 5/5 & 0/5 & 5/5 \\
004 exclude stale fact & 0/5 & 0/5 & 0/5 & 5/5 & 0/5 & 0/5 & 0/5 \\
\bottomrule
\end{tabular}
\caption{Deterministic v0.3 retrieval assertions passed / attempted across
five repetitions. A zero in a non-target arm is expected: only namespacing,
the gate, and the governed bundle exclude cross-domain distractors, while TTL
removes the stale fact. For nongated arms, retrieval is a deterministic
function of the manifest, configuration, and TF-IDF backend, so the five
repetitions re-observe an invariant result; gated rows remain repeated
observations of a model-mediated decision.}
\label{tab:retrieval}
\end{table}

\subsection{Attribution: what the cells support}

The single-control arms support only the claims their evidence layer can
bear. On the provenance probe, the tag and raw-fidelity arms each had five
machine-resolved clean outcomes. Semantic drift remains a response-layer tie:
all configurations resolved clean, while retrieval assertions distinguish the
candidate sets. Scope bleed separates at the response layer: naive resolved
one clean and four contaminated repetitions, whereas namespacing resolved five
clean and zero contaminated repetitions, with no flags in either cell. The
corresponding contamination proportions are 4/5 (exact 95\% Clopper--Pearson
interval [0.284, 0.995]) and 0/5 ([0.000, 0.522]); these intervals are
descriptive at validation scale, not population estimates. Their binomial
interpretation assumes independent trials with a fixed outcome probability;
repeated API calls do not establish that assumption or add scenario diversity.
Staleness remains
retrieval-layer evidence because its response-layer comparators are flagged.
The summarization probe has no utility oracle, so it cannot establish a
utility gain from raw fidelity.

The off-diagonal cells motivate two case-specific interpretations. First,
\emph{raw fidelity can preserve provenance in-band}: the raw store retains
the sentence ``the assistant suggested that Maria could try\ldots,'' and
the subject's responses were scored clean with no tags present. This is
consistent with loss of source information during summarization in this
scenario, but does not establish summarization as the principal cause of
provenance collapse generally. The two arms are alternative interventions
on the same case, not independent replications of that explanation.
Second, an implementation
confound we report rather than hide: the provenance tag as implemented
carries age and domain metadata (\texttt{[source: user | age: 200d |
domain: personal]}), so the provenance arm also resolved staleness because
the subject discounted the 200-day-old fact by its visible age. The tag
bundles three annotations; a stricter ablation would separate them, and
the full benchmark will.

\subsection{The gate amplified seeded recursion}
\label{sec:gatefinding}

In the frozen v0.2 single run, the relevance-gate arm was the only
configuration whose seeded-recursion verdict resolved contaminated, with
deterministic and judge passes in agreement. This exclusivity does not hold
in the v0.3 repeated audit: Table~\ref{tab:repeated} shows four contaminated
and one flagged gate outcome, and two contaminated and three flagged
provenance-arm outcomes. The other arms are entirely flagged on that probe,
so they cannot be treated as clean comparators or used to rank gate harm.
The v0.2 gate failure is traceable in
the gate's own logged decisions. Asked for the ``confirmed launch
date,'' the gate kept the contaminated write-back (``confirms Atlas
launch timing\ldots directly answers the query'') and discarded the
user-sourced ground truth as irrelevant, reasoning that it ``explicitly
states the launch date is tentative and not confirmed.'' The gate judged
relevance against the question's presupposition, stripped the
contradicting context, and handed the subject only the contamination;
the subject then asserted the confirmed launch as fact. Under naive
retrieval, with both memories present and no filter, the same subject
flagged the conflict and declined.

We call this \emph{presupposition capture}: query-aligned relevance is not
truth preservation, because a memory affirming a false premise can be more
responsive to the query than a memory that corrects it. The enforceable rail
is a retrieval contract: a probabilistic filter may never drop a memory that
contradicts a retained memory; contradiction sets are retained atomically.
The contract is implemented and unit-tested, but has not yet been evaluated
in a dedicated repeated run, so it is a safeguard proposal rather than an
efficacy result.

We scope this claim carefully: in this seeded-recursion validation
scenario, this relevance gate as implemented retained the contaminated
write-back and discarded the contradicting user context; generalizing to
relevance gates as a class requires the gate-family scenarios named in
future work. This suggests a structural risk when relevance is judged against a
presuppositional query: a relevance filter optimizes for answering the
question as asked, and a contaminated write-back is often the most
``relevant'' item in the store for precisely the question it contaminates.
The gate-family scenarios of the future-work contract
are intended to test this hypothesis. This case motivates the design
proposal in Section~\ref{sec:architecture}: the three
deterministic controls (namespacing, tags, decay) showed target-aligned
retrieval or resolved-response effects in this corpus without extra model
calls for enforcement; prompt-token and storage costs remain.
The v0.2 personalization-retention control resolved clean for six of
seven configurations, showing useful recall on that particular control,
not general preservation of utility. This is a design rationale, not a
general ordering of mitigation efficacy. The one
probabilistic control produced a recall false negative in the validation
run and discarded premise-denying context in the v0.2 seeded trace. These
observations motivate a targeted experiment, not a general claim of harm.

\subsection{Controls, retention, and remaining defects}
\label{sec:defects}

In the v0.2 single run, the anti-gaming control executed. With the over-broad pattern
from the prior run rebound to peanut-as-ingredient, the personalization
control resolved clean on six of seven configurations. All six
configurations with a machine-resolved retention result used relevant memory
correctly; the governed retention control remains flagged and is not counted
as a resolved success. This provides evidence against complete memory
suppression on this control in the six resolved configurations, but does not
measure broader utility or the contamination--retention trade-off.
The residual flag (governed) is a
negation case. The response mentioned a peanut ingredient in order to
warn against it, yet remains \texttt{needs\_human\_review} in the frozen
machine layer. The v0.2 defect note describes it as clean, but that note is
not a blinded, independent adjudication and is not used as a formal result.
It instead records the scoring defect: pattern-based scoring cannot see
negation or quotation, which is also the cause of the ten flagged artifacts
above. Deterministic patterns remain the first pass for auditability, but
assertion-versus-mention is precisely where the judge earns its place, and
the full benchmark will pair patterns with graded rubrics.

Two further honest observations. The natural-recursion compounding
factor is again null (\texttt{round1\_clean}) across all seven
configurations, so the seeded scenario now carries the mechanism test,
which is why it exists. And the borderline scope-bleed pattern
(\texttt{permit delays?}) fired under naive but not under three arms
whose prompts for that scenario are byte-identical to naive's:
temperature-0 sampling is not bit-reproducible across calls, so
borderline pattern hits are stochastic. The benchmark's reproducibility
claim attaches to the persisted artifacts, not to re-execution, and the
full benchmark averages repeated runs.

\section{A Contract-Level Mitigation Architecture}
\label{sec:architecture}

\paragraph{Declaration of interest.} The governed configuration instantiates
controls from the author's open COGNITION.md proposal. This paper therefore
reports a pilot control comparison, not an independent evaluation of
that proposal.

We use declarative constraints on writing, retrieval, injection, and
write-back as a design lens for the six mechanisms. The pilot does not
establish absence of such constraints as their shared or exclusive root cause.
Each baseline control in Table~\ref{tab:configs} is a contract-level answer
to a specific class: namespacing answers scope bleed; provenance tags answer
provenance collapse; decay by fact class answers temporal staleness; and
raw-fidelity storage preserves the information needed to assess
summarization loss. Relevance gating is not a baseline mitigation: it is an
optional residual-risk experiment that must obey the contradiction-pair rail
above before deployment.

The architectural principle is that these controls are \emph{contracts
external to the mutable store}, not behaviors the store is trusted to
exhibit. This is not a new principle; it is the reference monitor applied
to memory: a mediation layer that is tamperproof, always invoked, and
small enough to verify~\cite{anderson1972}, together with complete
mediation and the separation of policy from
mechanism~\cite{saltzer1975protection}. The harness literature reaches
the same conclusion for self-modifying code: when a system can edit its
own operating layer, safety controls must live outside the
loop~\cite{weng2026harness}. Besanson's SARC framework makes the analogous
case for tool-using agents: constraints should be first-class specifications
with explicit runtime enforcement sites and auditable traces~\cite{besanson2026sarc}.
Memory is the other self-editable surface.

The ablation motivates constraints, not an efficacy ordering. Namespacing,
provenance tags, TTL decay, and raw fidelity are deterministic in this harness
and require no additional model calls for their enforcement. They are not
cost-free: storage, metadata management, and prompt-token costs remain.
Their observed effects occur at
different evidence layers and in different classes. Relevance gating is a
probabilistic filter on the critical path. The repeated audit observed 135
gate calls over 100 gated retrievals, or 1.35 calls per gated retrieval; its
$k{=}4$ worst-case bound is 4{,}000 calls per thousand gated retrievals. No
batching is used. It is the only control that produced a recall false negative
in the validation run and active amplification of seeded recursion here
(Section~\ref{sec:gatefinding}). A deterministic, contract-enforced design
without a gate is therefore a reasonable baseline recommendation, not a
proven general mitigation ordering; a gate remains an unvalidated
residual-risk experiment, with logged decisions and mandatory contradiction
preservation. Relocating a probabilistic filter from the subject model into
the harness makes it auditable, but does not make it deterministic. This
aligns with the operating-system analogy in that literature: encapsulate
complex logic behind simple interfaces, and expect configuration and
interface protocols to standardize. The author's open specification
proposal, COGNITION.md~\cite{arana2026cognition}, is one such
standardization proposal for the memory subsystem specifically: scoping,
provenance, relevance gating, and decay expressed as first-class elements of
a declarative memory contract.

Production platforms are converging on the same controls from the
deployment side. This is a demand signal for memory lifecycle governance, not
evidence of correctness; only evaluation of the kind attempted here can
provide that. Anthropic's managed memory stores ship per-user and
per-team store scoping, read-only versus read-write access modes, and
immutable memory versions with audit trails~\cite{anthropic2026memory};
its Dreams feature performs asynchronous memory consolidation, motivated
explicitly by stores accumulating ``duplicates, contradictions, and stale
entries'' across sessions~\cite{anthropic2026dreams}. These capabilities address
memory lifecycle management; they are not evaluated in this pilot. We make no
comparative claim about their contamination rates or the completeness of their
safeguards. Consolidation summarizes and rewrites at store scale and may erase
provenance before retrieval. Its potential benefits and risks should be benchmarked
against the summarization-loss and provenance mechanisms of
Section~\ref{sec:taxonomy}.

\section{Discussion and Limitations}

\subsection{Validity statements}

Five limits bound what this study can claim, stated as validity
conditions rather than footnotes. (1)~Regex scoring detects
\emph{mention}, not assertion: quotation, negation, correction, and
refusal can match a forbidden phrase while being clean. In the v0.3
repeated audit, 52 of 350 rounds disagreed between deterministic and judge
scoring. Those rounds remain unresolved and excluded from machine-only
rates: the versioned adjudication queue assigns no human verdicts, so this
paper does not claim a judge-plus-human consensus for them. The correction
bundle records 0 adjudications and 0 two-adjudicator consensuses. (2)~The
provenance arm is confounded: its tag emits source, age, and domain
together, so its staleness result cannot be attributed to source
attribution alone. (3)~TF-IDF is a documented retrieval substitution; nothing here is evidence
about learned-embedding retrieval behavior. (4)~Temperature~0 reduces but
does not remove output variance, observed directly within this run on
byte-identical prompts. (5)~One subject
model cannot establish model-independent behavior, for compounding or
anything else. Rates from one probe per mechanism are mechanism
demonstrations with stated resolved counts, not population estimates. The
repeated audit's 660 API calls exceed the validation plan's ``few hundred
calls'' target; this is a cost limitation, not a reason to discard the
persisted artifacts. The review register identifies every flagged artifact by
repetition and hash.

\subsection{Future-work contract}

The next experiment is a preregistered full-benchmark study with at least
five scenarios per contamination class and paired relevant-memory controls,
repeated trials with per-repetition artifacts and variance reporting, at
least two subject models, and a learned-embedding retrieval backend reported
separately from TF-IDF. It will use pure-factor provenance arms,
production-default baselines (for example Mem0, Zep, and Letta), and a
declared utility oracle rather than treating retention as general utility.
It will also apply assertion-aware judge rubrics (asserted,
mentioned-but-rejected, quoted, hedged, or clean) and a preregistered,
blinded two-adjudicator protocol for every pending round, retaining only
independent matching verdicts as consensus. A dedicated gate-family suite
will measure precision, recall, and truth preservation under presupposition
traps, including the contradiction-pair rail. Together these measurements
will support a contamination-versus-retention frontier instead of the seven
mechanism-isolation points reported here.

The repository now includes a fail-closed preflight for this successor study:
it checks the planned 30 probes and 30 paired controls before authorization,
reports missing manifests, and estimates 240 subject calls for the two-model,
two-retrieval-condition matrix. At the time of writing, the semantic-drift
block supplies five probes and five paired controls; the remaining manifests,
gate-family suite, production-baseline adapters, and authorized execution are
not results and are not used by any claim in this paper.

\label{sec:discussion}

\paragraph{Scale.} The corpus is deliberately small: nine hand-authored
scenarios, one probe per class plus controls and the seeded variant,
designed to validate the pipeline and isolate mechanisms. The full
benchmark plan, versioned in the repository's roadmap and summarized
in the future-work contract above, extends the corpus to at least five
scenarios per class with paired controls, and extends the configuration
set with pure-factor arms; rates computed on single-probe classes take
values in $\{0, 1, \text{null}\}$ and are reported as rates only so the
formulas scale unchanged.

\paragraph{Severity proxy.} The compounding factor uses deterministic hit
counts as a coarse severity proxy; the full benchmark may adopt a graded
judge rubric.

\paragraph{Similarity substitution.} TF-IDF cosine similarity differs from
learned embeddings. A local ONNX learned-embedding condition is implemented
for the successor study, but it has not been run as part of the evidence
release. The hypothesis that learned embeddings may expose more semantic-drift
distractors therefore still requires a separately reported experiment; this
study makes no directional claim.

\paragraph{Single subject model.} Compounding is expected to be
model-dependent~\cite{zelikman2023stop}; the full benchmark evaluates
multiple subject models.

\paragraph{TTL classification.} Decay by fact class presupposes a
classifier that assigns the class at write time, and this paper does not
specify one. The default for unclassified memories (stable, 365 days) is
conservative for recall and anti-conservative for staleness. The axis
must be named because the two pull opposite ways. If an LLM assigns fact
classes, a probabilistic component enters the write path and class
misassignment becomes a failure mode not yet in the taxonomy; the full
benchmark must test it.

\paragraph{Toward the frontier.} Seven single-control arms are seven
points, not a curve. The prescriptive question is how much contamination is
bought down per unit of personalization surrendered. That trade-off is a
frontier, and mapping it (by sweeping control combinations, plotting
contamination rate against retention, locating the knee) is the successor
study this benchmark's full corpus is designed
for.

\paragraph{Reproducibility.} Temperature-0 sampling reduces but does not
guarantee bit-identical responses across executions of the same prompt.
Deterministic checks and aggregates can be recomputed from persisted
artifacts and recorded verdicts. LLM judge verdicts are retained as observations,
not guaranteed outputs of a new judge call. Hash inventories identify file
contents; they do not validate scoring correctness or independently reproduce
the study. Independent external reproduction and blinded human assessment
remain outstanding.

\paragraph{Future work.} Contamination in self-improving harness loops is
the natural extension: when a harness curates its own context from its own
trajectories, provenance collapse becomes a first-class risk of the
optimization loop itself. The experimental authority-inflation class
requires the multi-round, citation-count scenario design sketched in the
public taxonomy.

\section{Conclusion}

This pilot operationalizes six proposed mechanisms of benign memory
contamination in a small, auditable scenario suite. Within the tested cases,
namespacing and TTL changed targeted retrieval outcomes, while tags and raw
storage yielded clean machine verdicts on the provenance probe. A logged
seeded-recursion failure illustrates how relevance filtering can discard
context that denies a query's premise. Natural recursive compounding was
not observed, and the proposed contradiction-preservation safeguard has not
yet undergone a repeated efficacy evaluation.

One subject model, TF-IDF retrieval, limited scenario diversity, bundled
metadata, and unresolved scoring disagreements restrict these findings to
the evaluated conditions. They motivate further tests of memory contracts,
not a validated general mitigation architecture. Independent assessment of
scoring, broader scenarios with paired controls, and cross-model evaluation
are needed before stronger claims. The benchmark,
specifications, and validation scenarios are public at
\url{https://github.com/arananet/contam-bench}.

\bibliographystyle{plain}
\bibliography{references}

\end{document}
